% Pipeline figure: How Aldermen Influence the Development Process
% Use: % Pipeline figure: How Aldermen Influence the Development Process
% Use: % Pipeline figure: How Aldermen Influence the Development Process
% Use: % Pipeline figure: How Aldermen Influence the Development Process
% Use: \input{figures/pipeline_tikz.tex} inside a frame environment
% Requires: \usetikzlibrary{calc,positioning,arrows.meta,fit}
%
% 5 overlays:
%   <1> Pre-application (unobserved)
%   <2> + City Council zoning approval (only if a zoning change is needed) + Dalton & Tomich quote
%   <3> + Permit filed and reviewed + Dept of Buildings nightly reports quote
%   <4> + Permit issued + anonymous developer quote
%   <5> + observed-in-permit-data bracket and summary annotation
% Process: Chicago Zoning Ordinance 17-13-0300 (map amendments) and 17-13-0600 (planned developments): the Committee on
% Zoning, Landmarks and Building Standards holds a hearing and the full City Council decides; permits follow.

% --- Colors ---
\definecolor{pipeobs}{HTML}{2B7A78}
\definecolor{pipeunobs}{HTML}{8B8680}
\definecolor{pipequote}{HTML}{F5F0EB}
\definecolor{pipeaccent}{HTML}{C44D34}

\begin{tikzpicture}[
    remember picture,
    stage/.style={
        rounded corners=4pt,
        minimum height=1.35cm,
        minimum width=3.3cm,
        text width=3.0cm,
        align=center,
        font=\small,
        inner sep=6pt,
    },
    observed/.style={
        stage,
        draw=pipeobs,
        fill=pipeobs!8,
        line width=0.6pt,
    },
    unobserved/.style={
        stage,
        draw=pipeunobs,
        fill=pipeunobs!6,
        line width=0.6pt,
        dashed,
    },
    quotebox/.style={
        rounded corners=3pt,
        draw=black!15,
        fill=pipequote,
        text width=6.5cm,
        align=left,
        inner sep=8pt,
        font=\scriptsize,
    },
    arrw/.style={
        -{Stealth[length=5pt]},
        line width=0.7pt,
        pipeobs!70,
    },
    brack/.style={
        line width=0.4pt,
        black!30,
    },
]

% =============================================
% PIPELINE STAGES (top row) - fixed coordinates
% =============================================

% Stage 1: Pre-application (always visible)
\node[unobserved] (s1) at (0, 0) {
    \textbf{Pre-application}\\[1pt]
    {\scriptsize\textit{Meet the alderman;}}\\[-2pt]
    {\scriptsize\textit{community meetings}}
};

% Stage 2: City Council zoning approval
\visible<2->{
    \node[unobserved] (s2) at (4.1, 0) {
        \textbf{City Council}\\[1pt]
        {\scriptsize Zoning Committee}\\[-2pt]
        {\scriptsize hearing; Council vote}
    };
    \node[font=\tiny, text=pipeunobs, anchor=south] at (s2.north) {if a zoning change is needed};
    \draw[arrw] (s1.east) -- (s2.west);
}

% Stage 3: Permit filed and reviewed
\visible<3->{
    \node[observed] (s3) at (8.2, 0) {
        \textbf{Permit filed \& reviewed}\\[1pt]
        {\scriptsize Dept.\ of Buildings:}\\[-2pt]
        {\scriptsize plan review, revisions}
    };
    \draw[arrw] (s2.east) -- (s3.west);
}

% Stage 4: Permit issued
\visible<4->{
    \node[observed] (s4) at (12.3, 0) {
        \textbf{Permit issued}\\[1pt]
        {\scriptsize Construction,}\\[-2pt]
        {\scriptsize inspections, occupancy}
    };
    \draw[arrw] (s3.east) -- (s4.west);
}

% =============================================
% BRACKETS (overlay 5)
% =============================================

% Both brackets sit at one height, just below the tallest stage box (s3, whose title wraps to two lines).
\visible<5->{
    \coordinate (bracket) at ([yshift=-4pt]s3.south);

    \draw[brack, dashed] (s1.west |- bracket) -- (s2.east |- bracket);
    \draw[brack, dashed] (s1.west |- bracket) -- ++(0,3pt);
    \draw[brack, dashed] (s2.east |- bracket) -- ++(0,3pt);
    \coordinate (mid12) at ($(s1.south)!0.5!(s2.south)$);
    \node[font=\tiny, text=pipeunobs] at ([yshift=-12pt]mid12 |- bracket) {Not in permit data};

    \draw[brack] (s3.west |- bracket) -- (s4.east |- bracket);
    \draw[brack] (s3.west |- bracket) -- ++(0,3pt);
    \draw[brack] (s4.east |- bracket) -- ++(0,3pt);
    \coordinate (mid34) at ($(s3.south)!0.5!(s4.south)$);
    \node[font=\tiny, text=pipeobs] at ([yshift=-12pt]mid34 |- bracket)
        {Observable in permit data $\;\rightarrow\;$ proxy for $T(s)$};
}

% =============================================
% QUOTES AND NOTES (one per overlay, centered below)
% =============================================

\coordinate (qspot) at (6.15, -2.6);

% --- Overlay 1: pre-application (deterrence) ---
\only<1>{
    \node[font=\scriptsize, text=black!70, text width=7cm, align=left, anchor=north] (q1) at (qspot) {
        {\bfseries\scriptsize\textcolor{pipeaccent}{\textsc{Deterrence}}}\\[4pt]
        Developers meet the ward's alderman and neighbors before filing. Projects discouraged here are never proposed.
    };
}

% --- Overlay 2: City Council (aldermanic deference on zoning) ---
\only<2>{
    \node[quotebox, anchor=north] (q2) at (qspot) {
        {\bfseries\scriptsize\textcolor{pipeaccent}{\textsc{Deference}}}\\[4pt]
        ``A thumb up or thumb down from an alderman can mean all the difference to even the most code-compliant zoning application.''\\[4pt]
        \hfill{\tiny\textcolor{black!45}{--- Dalton \& Tomich, land use attorneys}}
    };
    \draw[brack, dashed] (s2.south) -- ([yshift=4pt]q2.north -| s2.south);
}

% --- Overlay 3: Dept of Buildings (monitoring) ---
\only<3>{
    \node[quotebox, anchor=north] (q3) at (qspot) {
        {\bfseries\scriptsize\textcolor{pipeaccent}{\textsc{Monitoring}}}\\[4pt]
        ``As a matter of courtesy, aldermen continue to receive nightly e-mailed reports of building permits applied for and issued in their respective wards.''\\[4pt]
        \hfill{\tiny\textcolor{black!45}{--- City of Chicago, Dept.\ of Buildings}}
    };
    \draw[brack, dashed] (s3.south) -- ([yshift=4pt]q3.north -| s3.south);
}

% --- Overlay 4: Anonymous developer (direct intervention) ---
\only<4>{
    \node[quotebox, anchor=north] (q4) at (qspot) {
        {\bfseries\scriptsize\textcolor{pipeaccent}{\textsc{Direct intervention}}}\\[4pt]
        ``If the alderman doesn't want it to happen, it won't happen. All it takes is one call to a city inspector, saying `Don't give him his occupancy permit.'\,''\\[4pt]
        \hfill{\tiny\textcolor{black!45}{--- Anonymous developer}}
    };
    \draw[brack, dashed] (s4.south) -- ([yshift=4pt]q4.north -| s4.south);
}

% --- Overlay 5: Summary ---
\only<5>{
    \node[font=\scriptsize, text=black!70, text width=7cm, align=left, anchor=north] at ([yshift=-0.35cm]qspot) {
        \textbf{My measure}: Alderman-specific component of permit processing time captures frictions in the observable stages.\\[3pt]
        \textit{Coverage}: Projects deterred before application, or stopped at the Council, never enter the permit data.
    };
}

\end{tikzpicture}
 inside a frame environment
% Requires: \usetikzlibrary{calc,positioning,arrows.meta,fit}
%
% 5 overlays:
%   <1> Pre-application (unobserved)
%   <2> + City Council zoning approval (only if a zoning change is needed) + Dalton & Tomich quote
%   <3> + Permit filed and reviewed + Dept of Buildings nightly reports quote
%   <4> + Permit issued + anonymous developer quote
%   <5> + observed-in-permit-data bracket and summary annotation
% Process: Chicago Zoning Ordinance 17-13-0300 (map amendments) and 17-13-0600 (planned developments): the Committee on
% Zoning, Landmarks and Building Standards holds a hearing and the full City Council decides; permits follow.

% --- Colors ---
\definecolor{pipeobs}{HTML}{2B7A78}
\definecolor{pipeunobs}{HTML}{8B8680}
\definecolor{pipequote}{HTML}{F5F0EB}
\definecolor{pipeaccent}{HTML}{C44D34}

\begin{tikzpicture}[
    remember picture,
    stage/.style={
        rounded corners=4pt,
        minimum height=1.35cm,
        minimum width=3.3cm,
        text width=3.0cm,
        align=center,
        font=\small,
        inner sep=6pt,
    },
    observed/.style={
        stage,
        draw=pipeobs,
        fill=pipeobs!8,
        line width=0.6pt,
    },
    unobserved/.style={
        stage,
        draw=pipeunobs,
        fill=pipeunobs!6,
        line width=0.6pt,
        dashed,
    },
    quotebox/.style={
        rounded corners=3pt,
        draw=black!15,
        fill=pipequote,
        text width=6.5cm,
        align=left,
        inner sep=8pt,
        font=\scriptsize,
    },
    arrw/.style={
        -{Stealth[length=5pt]},
        line width=0.7pt,
        pipeobs!70,
    },
    brack/.style={
        line width=0.4pt,
        black!30,
    },
]

% =============================================
% PIPELINE STAGES (top row) - fixed coordinates
% =============================================

% Stage 1: Pre-application (always visible)
\node[unobserved] (s1) at (0, 0) {
    \textbf{Pre-application}\\[1pt]
    {\scriptsize\textit{Meet the alderman;}}\\[-2pt]
    {\scriptsize\textit{community meetings}}
};

% Stage 2: City Council zoning approval
\visible<2->{
    \node[unobserved] (s2) at (4.1, 0) {
        \textbf{City Council}\\[1pt]
        {\scriptsize Zoning Committee}\\[-2pt]
        {\scriptsize hearing; Council vote}
    };
    \node[font=\tiny, text=pipeunobs, anchor=south] at (s2.north) {if a zoning change is needed};
    \draw[arrw] (s1.east) -- (s2.west);
}

% Stage 3: Permit filed and reviewed
\visible<3->{
    \node[observed] (s3) at (8.2, 0) {
        \textbf{Permit filed \& reviewed}\\[1pt]
        {\scriptsize Dept.\ of Buildings:}\\[-2pt]
        {\scriptsize plan review, revisions}
    };
    \draw[arrw] (s2.east) -- (s3.west);
}

% Stage 4: Permit issued
\visible<4->{
    \node[observed] (s4) at (12.3, 0) {
        \textbf{Permit issued}\\[1pt]
        {\scriptsize Construction,}\\[-2pt]
        {\scriptsize inspections, occupancy}
    };
    \draw[arrw] (s3.east) -- (s4.west);
}

% =============================================
% BRACKETS (overlay 5)
% =============================================

% Both brackets sit at one height, just below the tallest stage box (s3, whose title wraps to two lines).
\visible<5->{
    \coordinate (bracket) at ([yshift=-4pt]s3.south);

    \draw[brack, dashed] (s1.west |- bracket) -- (s2.east |- bracket);
    \draw[brack, dashed] (s1.west |- bracket) -- ++(0,3pt);
    \draw[brack, dashed] (s2.east |- bracket) -- ++(0,3pt);
    \coordinate (mid12) at ($(s1.south)!0.5!(s2.south)$);
    \node[font=\tiny, text=pipeunobs] at ([yshift=-12pt]mid12 |- bracket) {Not in permit data};

    \draw[brack] (s3.west |- bracket) -- (s4.east |- bracket);
    \draw[brack] (s3.west |- bracket) -- ++(0,3pt);
    \draw[brack] (s4.east |- bracket) -- ++(0,3pt);
    \coordinate (mid34) at ($(s3.south)!0.5!(s4.south)$);
    \node[font=\tiny, text=pipeobs] at ([yshift=-12pt]mid34 |- bracket)
        {Observable in permit data $\;\rightarrow\;$ proxy for $T(s)$};
}

% =============================================
% QUOTES AND NOTES (one per overlay, centered below)
% =============================================

\coordinate (qspot) at (6.15, -2.6);

% --- Overlay 1: pre-application (deterrence) ---
\only<1>{
    \node[font=\scriptsize, text=black!70, text width=7cm, align=left, anchor=north] (q1) at (qspot) {
        {\bfseries\scriptsize\textcolor{pipeaccent}{\textsc{Deterrence}}}\\[4pt]
        Developers meet the ward's alderman and neighbors before filing. Projects discouraged here are never proposed.
    };
}

% --- Overlay 2: City Council (aldermanic deference on zoning) ---
\only<2>{
    \node[quotebox, anchor=north] (q2) at (qspot) {
        {\bfseries\scriptsize\textcolor{pipeaccent}{\textsc{Deference}}}\\[4pt]
        ``A thumb up or thumb down from an alderman can mean all the difference to even the most code-compliant zoning application.''\\[4pt]
        \hfill{\tiny\textcolor{black!45}{--- Dalton \& Tomich, land use attorneys}}
    };
    \draw[brack, dashed] (s2.south) -- ([yshift=4pt]q2.north -| s2.south);
}

% --- Overlay 3: Dept of Buildings (monitoring) ---
\only<3>{
    \node[quotebox, anchor=north] (q3) at (qspot) {
        {\bfseries\scriptsize\textcolor{pipeaccent}{\textsc{Monitoring}}}\\[4pt]
        ``As a matter of courtesy, aldermen continue to receive nightly e-mailed reports of building permits applied for and issued in their respective wards.''\\[4pt]
        \hfill{\tiny\textcolor{black!45}{--- City of Chicago, Dept.\ of Buildings}}
    };
    \draw[brack, dashed] (s3.south) -- ([yshift=4pt]q3.north -| s3.south);
}

% --- Overlay 4: Anonymous developer (direct intervention) ---
\only<4>{
    \node[quotebox, anchor=north] (q4) at (qspot) {
        {\bfseries\scriptsize\textcolor{pipeaccent}{\textsc{Direct intervention}}}\\[4pt]
        ``If the alderman doesn't want it to happen, it won't happen. All it takes is one call to a city inspector, saying `Don't give him his occupancy permit.'\,''\\[4pt]
        \hfill{\tiny\textcolor{black!45}{--- Anonymous developer}}
    };
    \draw[brack, dashed] (s4.south) -- ([yshift=4pt]q4.north -| s4.south);
}

% --- Overlay 5: Summary ---
\only<5>{
    \node[font=\scriptsize, text=black!70, text width=7cm, align=left, anchor=north] at ([yshift=-0.35cm]qspot) {
        \textbf{My measure}: Alderman-specific component of permit processing time captures frictions in the observable stages.\\[3pt]
        \textit{Coverage}: Projects deterred before application, or stopped at the Council, never enter the permit data.
    };
}

\end{tikzpicture}
 inside a frame environment
% Requires: \usetikzlibrary{calc,positioning,arrows.meta,fit}
%
% 5 overlays:
%   <1> Pre-application (unobserved)
%   <2> + City Council zoning approval (only if a zoning change is needed) + Dalton & Tomich quote
%   <3> + Permit filed and reviewed + Dept of Buildings nightly reports quote
%   <4> + Permit issued + anonymous developer quote
%   <5> + observed-in-permit-data bracket and summary annotation
% Process: Chicago Zoning Ordinance 17-13-0300 (map amendments) and 17-13-0600 (planned developments): the Committee on
% Zoning, Landmarks and Building Standards holds a hearing and the full City Council decides; permits follow.

% --- Colors ---
\definecolor{pipeobs}{HTML}{2B7A78}
\definecolor{pipeunobs}{HTML}{8B8680}
\definecolor{pipequote}{HTML}{F5F0EB}
\definecolor{pipeaccent}{HTML}{C44D34}

\begin{tikzpicture}[
    remember picture,
    stage/.style={
        rounded corners=4pt,
        minimum height=1.35cm,
        minimum width=3.3cm,
        text width=3.0cm,
        align=center,
        font=\small,
        inner sep=6pt,
    },
    observed/.style={
        stage,
        draw=pipeobs,
        fill=pipeobs!8,
        line width=0.6pt,
    },
    unobserved/.style={
        stage,
        draw=pipeunobs,
        fill=pipeunobs!6,
        line width=0.6pt,
        dashed,
    },
    quotebox/.style={
        rounded corners=3pt,
        draw=black!15,
        fill=pipequote,
        text width=6.5cm,
        align=left,
        inner sep=8pt,
        font=\scriptsize,
    },
    arrw/.style={
        -{Stealth[length=5pt]},
        line width=0.7pt,
        pipeobs!70,
    },
    brack/.style={
        line width=0.4pt,
        black!30,
    },
]

% =============================================
% PIPELINE STAGES (top row) - fixed coordinates
% =============================================

% Stage 1: Pre-application (always visible)
\node[unobserved] (s1) at (0, 0) {
    \textbf{Pre-application}\\[1pt]
    {\scriptsize\textit{Meet the alderman;}}\\[-2pt]
    {\scriptsize\textit{community meetings}}
};

% Stage 2: City Council zoning approval
\visible<2->{
    \node[unobserved] (s2) at (4.1, 0) {
        \textbf{City Council}\\[1pt]
        {\scriptsize Zoning Committee}\\[-2pt]
        {\scriptsize hearing; Council vote}
    };
    \node[font=\tiny, text=pipeunobs, anchor=south] at (s2.north) {if a zoning change is needed};
    \draw[arrw] (s1.east) -- (s2.west);
}

% Stage 3: Permit filed and reviewed
\visible<3->{
    \node[observed] (s3) at (8.2, 0) {
        \textbf{Permit filed \& reviewed}\\[1pt]
        {\scriptsize Dept.\ of Buildings:}\\[-2pt]
        {\scriptsize plan review, revisions}
    };
    \draw[arrw] (s2.east) -- (s3.west);
}

% Stage 4: Permit issued
\visible<4->{
    \node[observed] (s4) at (12.3, 0) {
        \textbf{Permit issued}\\[1pt]
        {\scriptsize Construction,}\\[-2pt]
        {\scriptsize inspections, occupancy}
    };
    \draw[arrw] (s3.east) -- (s4.west);
}

% =============================================
% BRACKETS (overlay 5)
% =============================================

% Both brackets sit at one height, just below the tallest stage box (s3, whose title wraps to two lines).
\visible<5->{
    \coordinate (bracket) at ([yshift=-4pt]s3.south);

    \draw[brack, dashed] (s1.west |- bracket) -- (s2.east |- bracket);
    \draw[brack, dashed] (s1.west |- bracket) -- ++(0,3pt);
    \draw[brack, dashed] (s2.east |- bracket) -- ++(0,3pt);
    \coordinate (mid12) at ($(s1.south)!0.5!(s2.south)$);
    \node[font=\tiny, text=pipeunobs] at ([yshift=-12pt]mid12 |- bracket) {Not in permit data};

    \draw[brack] (s3.west |- bracket) -- (s4.east |- bracket);
    \draw[brack] (s3.west |- bracket) -- ++(0,3pt);
    \draw[brack] (s4.east |- bracket) -- ++(0,3pt);
    \coordinate (mid34) at ($(s3.south)!0.5!(s4.south)$);
    \node[font=\tiny, text=pipeobs] at ([yshift=-12pt]mid34 |- bracket)
        {Observable in permit data $\;\rightarrow\;$ proxy for $T(s)$};
}

% =============================================
% QUOTES AND NOTES (one per overlay, centered below)
% =============================================

\coordinate (qspot) at (6.15, -2.6);

% --- Overlay 1: pre-application (deterrence) ---
\only<1>{
    \node[font=\scriptsize, text=black!70, text width=7cm, align=left, anchor=north] (q1) at (qspot) {
        {\bfseries\scriptsize\textcolor{pipeaccent}{\textsc{Deterrence}}}\\[4pt]
        Developers meet the ward's alderman and neighbors before filing. Projects discouraged here are never proposed.
    };
}

% --- Overlay 2: City Council (aldermanic deference on zoning) ---
\only<2>{
    \node[quotebox, anchor=north] (q2) at (qspot) {
        {\bfseries\scriptsize\textcolor{pipeaccent}{\textsc{Deference}}}\\[4pt]
        ``A thumb up or thumb down from an alderman can mean all the difference to even the most code-compliant zoning application.''\\[4pt]
        \hfill{\tiny\textcolor{black!45}{--- Dalton \& Tomich, land use attorneys}}
    };
    \draw[brack, dashed] (s2.south) -- ([yshift=4pt]q2.north -| s2.south);
}

% --- Overlay 3: Dept of Buildings (monitoring) ---
\only<3>{
    \node[quotebox, anchor=north] (q3) at (qspot) {
        {\bfseries\scriptsize\textcolor{pipeaccent}{\textsc{Monitoring}}}\\[4pt]
        ``As a matter of courtesy, aldermen continue to receive nightly e-mailed reports of building permits applied for and issued in their respective wards.''\\[4pt]
        \hfill{\tiny\textcolor{black!45}{--- City of Chicago, Dept.\ of Buildings}}
    };
    \draw[brack, dashed] (s3.south) -- ([yshift=4pt]q3.north -| s3.south);
}

% --- Overlay 4: Anonymous developer (direct intervention) ---
\only<4>{
    \node[quotebox, anchor=north] (q4) at (qspot) {
        {\bfseries\scriptsize\textcolor{pipeaccent}{\textsc{Direct intervention}}}\\[4pt]
        ``If the alderman doesn't want it to happen, it won't happen. All it takes is one call to a city inspector, saying `Don't give him his occupancy permit.'\,''\\[4pt]
        \hfill{\tiny\textcolor{black!45}{--- Anonymous developer}}
    };
    \draw[brack, dashed] (s4.south) -- ([yshift=4pt]q4.north -| s4.south);
}

% --- Overlay 5: Summary ---
\only<5>{
    \node[font=\scriptsize, text=black!70, text width=7cm, align=left, anchor=north] at ([yshift=-0.35cm]qspot) {
        \textbf{My measure}: Alderman-specific component of permit processing time captures frictions in the observable stages.\\[3pt]
        \textit{Coverage}: Projects deterred before application, or stopped at the Council, never enter the permit data.
    };
}

\end{tikzpicture}
 inside a frame environment
% Requires: \usetikzlibrary{calc,positioning,arrows.meta,fit}
%
% 5 overlays:
%   <1> Pre-application (unobserved)
%   <2> + City Council zoning approval (only if a zoning change is needed) + Dalton & Tomich quote
%   <3> + Permit filed and reviewed + Dept of Buildings nightly reports quote
%   <4> + Permit issued + anonymous developer quote
%   <5> + observed-in-permit-data bracket and summary annotation
% Process: Chicago Zoning Ordinance 17-13-0300 (map amendments) and 17-13-0600 (planned developments): the Committee on
% Zoning, Landmarks and Building Standards holds a hearing and the full City Council decides; permits follow.

% --- Colors ---
\definecolor{pipeobs}{HTML}{2B7A78}
\definecolor{pipeunobs}{HTML}{8B8680}
\definecolor{pipequote}{HTML}{F5F0EB}
\definecolor{pipeaccent}{HTML}{C44D34}

\begin{tikzpicture}[
    remember picture,
    stage/.style={
        rounded corners=4pt,
        minimum height=1.35cm,
        minimum width=3.3cm,
        text width=3.0cm,
        align=center,
        font=\small,
        inner sep=6pt,
    },
    observed/.style={
        stage,
        draw=pipeobs,
        fill=pipeobs!8,
        line width=0.6pt,
    },
    unobserved/.style={
        stage,
        draw=pipeunobs,
        fill=pipeunobs!6,
        line width=0.6pt,
        dashed,
    },
    quotebox/.style={
        rounded corners=3pt,
        draw=black!15,
        fill=pipequote,
        text width=6.5cm,
        align=left,
        inner sep=8pt,
        font=\scriptsize,
    },
    arrw/.style={
        -{Stealth[length=5pt]},
        line width=0.7pt,
        pipeobs!70,
    },
    brack/.style={
        line width=0.4pt,
        black!30,
    },
]

% =============================================
% PIPELINE STAGES (top row) - fixed coordinates
% =============================================

% Stage 1: Pre-application (always visible)
\node[unobserved] (s1) at (0, 0) {
    \textbf{Pre-application}\\[1pt]
    {\scriptsize\textit{Meet the alderman;}}\\[-2pt]
    {\scriptsize\textit{community meetings}}
};

% Stage 2: City Council zoning approval
\visible<2->{
    \node[unobserved] (s2) at (4.1, 0) {
        \textbf{City Council}\\[1pt]
        {\scriptsize Zoning Committee}\\[-2pt]
        {\scriptsize hearing; Council vote}
    };
    \node[font=\tiny, text=pipeunobs, anchor=south] at (s2.north) {if a zoning change is needed};
    \draw[arrw] (s1.east) -- (s2.west);
}

% Stage 3: Permit filed and reviewed
\visible<3->{
    \node[observed] (s3) at (8.2, 0) {
        \textbf{Permit filed \& reviewed}\\[1pt]
        {\scriptsize Dept.\ of Buildings:}\\[-2pt]
        {\scriptsize plan review, revisions}
    };
    \draw[arrw] (s2.east) -- (s3.west);
}

% Stage 4: Permit issued
\visible<4->{
    \node[observed] (s4) at (12.3, 0) {
        \textbf{Permit issued}\\[1pt]
        {\scriptsize Construction,}\\[-2pt]
        {\scriptsize inspections, occupancy}
    };
    \draw[arrw] (s3.east) -- (s4.west);
}

% =============================================
% BRACKETS (overlay 5)
% =============================================

% Both brackets sit at one height, just below the tallest stage box (s3, whose title wraps to two lines).
\visible<5->{
    \coordinate (bracket) at ([yshift=-4pt]s3.south);

    \draw[brack, dashed] (s1.west |- bracket) -- (s2.east |- bracket);
    \draw[brack, dashed] (s1.west |- bracket) -- ++(0,3pt);
    \draw[brack, dashed] (s2.east |- bracket) -- ++(0,3pt);
    \coordinate (mid12) at ($(s1.south)!0.5!(s2.south)$);
    \node[font=\tiny, text=pipeunobs] at ([yshift=-12pt]mid12 |- bracket) {Not in permit data};

    \draw[brack] (s3.west |- bracket) -- (s4.east |- bracket);
    \draw[brack] (s3.west |- bracket) -- ++(0,3pt);
    \draw[brack] (s4.east |- bracket) -- ++(0,3pt);
    \coordinate (mid34) at ($(s3.south)!0.5!(s4.south)$);
    \node[font=\tiny, text=pipeobs] at ([yshift=-12pt]mid34 |- bracket)
        {Observable in permit data $\;\rightarrow\;$ proxy for $T(s)$};
}

% =============================================
% QUOTES AND NOTES (one per overlay, centered below)
% =============================================

\coordinate (qspot) at (6.15, -2.6);

% --- Overlay 1: pre-application (deterrence) ---
\only<1>{
    \node[font=\scriptsize, text=black!70, text width=7cm, align=left, anchor=north] (q1) at (qspot) {
        {\bfseries\scriptsize\textcolor{pipeaccent}{\textsc{Deterrence}}}\\[4pt]
        Developers meet the ward's alderman and neighbors before filing. Projects discouraged here are never proposed.
    };
}

% --- Overlay 2: City Council (aldermanic deference on zoning) ---
\only<2>{
    \node[quotebox, anchor=north] (q2) at (qspot) {
        {\bfseries\scriptsize\textcolor{pipeaccent}{\textsc{Deference}}}\\[4pt]
        ``A thumb up or thumb down from an alderman can mean all the difference to even the most code-compliant zoning application.''\\[4pt]
        \hfill{\tiny\textcolor{black!45}{--- Dalton \& Tomich, land use attorneys}}
    };
    \draw[brack, dashed] (s2.south) -- ([yshift=4pt]q2.north -| s2.south);
}

% --- Overlay 3: Dept of Buildings (monitoring) ---
\only<3>{
    \node[quotebox, anchor=north] (q3) at (qspot) {
        {\bfseries\scriptsize\textcolor{pipeaccent}{\textsc{Monitoring}}}\\[4pt]
        ``As a matter of courtesy, aldermen continue to receive nightly e-mailed reports of building permits applied for and issued in their respective wards.''\\[4pt]
        \hfill{\tiny\textcolor{black!45}{--- City of Chicago, Dept.\ of Buildings}}
    };
    \draw[brack, dashed] (s3.south) -- ([yshift=4pt]q3.north -| s3.south);
}

% --- Overlay 4: Anonymous developer (direct intervention) ---
\only<4>{
    \node[quotebox, anchor=north] (q4) at (qspot) {
        {\bfseries\scriptsize\textcolor{pipeaccent}{\textsc{Direct intervention}}}\\[4pt]
        ``If the alderman doesn't want it to happen, it won't happen. All it takes is one call to a city inspector, saying `Don't give him his occupancy permit.'\,''\\[4pt]
        \hfill{\tiny\textcolor{black!45}{--- Anonymous developer}}
    };
    \draw[brack, dashed] (s4.south) -- ([yshift=4pt]q4.north -| s4.south);
}

% --- Overlay 5: Summary ---
\only<5>{
    \node[font=\scriptsize, text=black!70, text width=7cm, align=left, anchor=north] at ([yshift=-0.35cm]qspot) {
        \textbf{My measure}: Alderman-specific component of permit processing time captures frictions in the observable stages.\\[3pt]
        \textit{Coverage}: Projects deterred before application, or stopped at the Council, never enter the permit data.
    };
}

\end{tikzpicture}
